\appendix
\ThesisPageBreak
\phantomsection
\addcontentsline{toc}{chapter}{\AppendixLabel}
\chapter*{\AppendixLabel}

% Penomoran materi lampiran dipisahkan dari penomoran bab utama agar tabel dan
% gambar pendukung tidak melanjutkan nomor Bab IV: Tabel L.1, Gambar L.1, dan
% seterusnya.
\setcounter{table}{0}
\setcounter{figure}{0}
\renewcommand{\thetable}{L.\arabic{table}}
\renewcommand{\thefigure}{L.\arabic{figure}}

\TemplateTodo{Gunakan lampiran untuk materi pendukung yang terlalu panjang jika ditempatkan di bab utama, misalnya instrumen penelitian, konfigurasi eksperimen, tabel hasil lengkap, bukti pengujian, dokumen etik, atau lembar validasi data.}

Tuliskan satu paragraf pengantar yang menyebutkan isi seluruh lampiran. Setiap lampiran harus dirujuk dari bab utama; jangan memasukkan materi yang tidak pernah disebutkan dalam naskah.

% Format lampiran:
%   - \AppendixSection{Lampiran 1. Judul}, Lampiran 2, dan seterusnya: lampiran
%     diberi nomor angka, bukan huruf. Lampiran kedua dan seterusnya otomatis
%     dimulai di halaman baru.
%   - \AppendixSubsection{Judul} tanpa nomor.
%   - Keduanya tercatat di DAFTAR LAMPIRAN, bukan di DAFTAR ISI.
%   - \label tepat setelah keduanya dapat dirujuk dari bab dengan \nameref,
%     yang mencetak judulnya: "Lampiran \nameref{sec:lampiran-contoh}".

\AppendixSection{Lampiran 1. Contoh Lampiran}\label{sec:lampiran-contoh}
Tuliskan isi lampiran di sini. Jika lampiran berupa kode, data, atau gambar, pastikan informasi sensitif sudah dihapus sebelum repositori dipublikasikan.

\AppendixSubsection{Contoh Subbagian Lampiran}\label{sec:lampiran-contoh-sub}
Rincian pendukung ditunjukkan pada Tabel \ref{tab:lampiran-contoh}.

\begin{longtable}{|l|c|c|}
\caption{Contoh tabel lampiran.}\label{tab:lampiran-contoh}\\
\hline
Konfigurasi & Metrik 1 & Metrik 2 \\
\hline
\endfirsthead
\hline
Konfigurasi & Metrik 1 & Metrik 2 \\
\hline
\endhead
\hline
\endlastfoot
Konfigurasi A & 0,00 & 0,00 \\
\hline
Konfigurasi B & 0,00 & 0,00 \\
\end{longtable}

\AppendixSection{Lampiran 2. Dokumen Pendukung}\label{sec:lampiran-dokumen}
Pindaian dokumen, misalnya surat keterangan etik atau lembar validasi data, disisipkan sebagai gambar berbingkai.

\begin{figure}[H]
\centering
\BorderedGraphics[width=0.6\linewidth]{example-figure.pdf}
\caption{Contoh pindaian dokumen pendukung.}
\label{fig:lampiran-dokumen}
\end{figure}
